The fixed-end eigenvalue analysis follows the course notes on string natural modes, PDF pp. 145--149.

\textbf{1--4. Boundary conditions.} Both pinned ends have zero transverse displacement. Using the solution from A.1,
\[
u(0,t)=u(L,t)=0,\qquad
U(x)=C_1\sin(\gamma x)+C_2\cos(\gamma x).
\]
The left boundary gives $C_2=0$, and the right boundary then requires $C_1\sin(\gamma L)=0$.

\textbf{5--9. Eigenvalues and wavenumbers.} $C_1=0$ produces zero displacement everywhere and is the trivial solution. A nonzero amplitude requires $\sin(\gamma L)=0$. Introducing $z=\gamma L$ gives
\[
\sin z=0,\qquad \boxed{z_j=j\pi,\quad
\gamma_j=\frac{j\pi}{L},\quad j=1,2,3,\ldots.}
\]
The zero root yields the zero shape under the fixed-end conditions and is excluded.

\textbf{10--14. Natural frequencies.} Since $\gamma=\omega/c$, the natural angular and cyclic frequencies are
\[
\boxed{\omega_j=\frac{j\pi c}{L},\qquad f_j=\frac{\omega_j}{2\pi}=\frac{jc}{2L}.}
\]
The fundamental is the lowest natural frequency $f_1$. The overtones are the higher frequencies $f_2,f_3,\ldots$.

\textbf{15--17. Wavelengths.} With $\gamma=2\pi/\lambda$, the allowable wavelengths satisfy
\[
\boxed{\lambda_j=\frac{2L}{j},\qquad L=j\frac{\lambda_j}{2}.}
\]
An integer number of half-wavelengths fits between the fixed ends.

\textbf{18--21. Modes and orthogonality.} Setting $C_2=0$ and choosing $C_1=1$ gives the unit-amplitude mode shape
\[
U_j(x)=\sin\left(\frac{j\pi x}{L}\right).
\]
These unit-amplitude shapes are orthogonal, with $\int_0^L U_jU_k\,dx=(L/2)\delta_{jk}$. To obtain orthonormal shapes under the unweighted inner product, set
\[
\phi_j(x)=\sqrt{\frac{2}{L}}\sin\left(\frac{j\pi x}{L}\right),
\qquad \boxed{\int_0^L\phi_j(x)\phi_k(x)\,dx=\delta_{jk}.}
\]
